\documentclass[11pt]{article}
\usepackage{amsmath,amssymb}
\usepackage{geometry}
\usepackage{parskip}
\usepackage{hyperref}
\geometry{margin=1.1in}

\title{Biased Collective-Coordinate Potential\\for Hole-Size Measurement}
\author{CCO-projects --- \texttt{biased\_optimization} branch}
\date{}

\begin{document}
\maketitle
\tableofcontents
\bigskip

% ------------------------------------------------------------------
\section{Original collective-coordinate potential}
% ------------------------------------------------------------------

The standard stealthy collective-coordinate (CC) potential is
\begin{equation}
    \Phi_{\mathrm{CC}}(\{\mathbf{r}_i\})
    = \frac{1}{V}
      \sum_{\mathbf{k}\in\mathcal{K}}
      \tilde{v}(\mathbf{k})\,|\rho(\mathbf{k})|^{2},
    \label{eq:cc}
\end{equation}
where
\begin{align}
    \rho(\mathbf{k}) &= \sum_{j=1}^{N} e^{i\mathbf{k}\cdot\mathbf{r}_j}
    \quad\text{(structure-factor amplitude)},\\[4pt]
    \mathcal{K} &= \bigcup_{n=1}^{M}
        \bigl\{\mathbf{k} : k_1^{(n)} < |\mathbf{k}| \le k_1^{(n)}+\delta^{(n)}\bigr\}
    \quad\text{(excluded $k$-space region)},
\end{align}
$V$ is the simulation-box volume, and $\tilde{v}(\mathbf{k})=1$ for the
flat (default) stealthy weight.
Minimising $\Phi_{\mathrm{CC}}$ to zero enforces $S(\mathbf{k})=0$ for
all $\mathbf{k}\in\mathcal{K}$, the defining condition of stealthy
hyperuniformity.

% ------------------------------------------------------------------
\section{Goal: maximum hole radius at the box centre}
% ------------------------------------------------------------------

We wish to measure the maximum radius $R_c$ of a spherical void that can
be placed at the box centre
\begin{equation}
    \mathbf{r}_c = \tfrac{1}{2}(\mathbf{a}_1 + \mathbf{a}_2 + \cdots + \mathbf{a}_d),
\end{equation}
where $\mathbf{a}_i$ are the lattice vectors, while the configuration
simultaneously satisfies the stealthy constraint
$\Phi_{\mathrm{CC}}\approx 0$.

% ------------------------------------------------------------------
\section{Exclusion field added to the potential}
% ------------------------------------------------------------------

To enforce a void of probe radius $R_f$ around $\mathbf{r}_c$, we add a
radial exclusion field that diverges as any particle approaches the
centre:
\begin{equation}
    \phi_{\mathrm{ex}}(\{\mathbf{r}_i\};\,R_f)
    = \sum_{\substack{i=1\\r_{ic}<R_f}}^{N}
      \left(\frac{R_f}{r_{ic}} - 1\right),
    \qquad
    r_{ic} = |\mathbf{r}_i - \mathbf{r}_c|.
    \label{eq:phiex}
\end{equation}
Key properties of $\phi_{\mathrm{ex}}$:
\begin{itemize}
    \item It is \emph{zero} at $r_{ic}=R_f$ (particle on the boundary).
    \item It \emph{diverges} as $r_{ic}\to 0$ (particle at the centre).
    \item It is \emph{non-negative} for all $r_{ic}>0$.
    \item Only particles strictly inside the probe sphere ($r_{ic}<R_f$)
          contribute.
\end{itemize}

% ------------------------------------------------------------------
\section{Biased potential}
% ------------------------------------------------------------------

The full biased potential used during minimisation is
\begin{equation}
    \boxed{
    \Phi_{\mathrm{biased}}(\{\mathbf{r}_i\};\,R_f)
    = \underbrace{\frac{1}{V}
      \sum_{\mathbf{k}\in\mathcal{K}}
      \tilde{v}(\mathbf{k})\,|\rho(\mathbf{k})|^{2}}_{\Phi_{\mathrm{CC}}}
    +
    \underbrace{\sum_{\substack{i\\r_{ic}<R_f}}
      \!\!\left(\frac{R_f}{r_{ic}}-1\right)}_{\phi_{\mathrm{ex}}}.
    }
    \label{eq:biased}
\end{equation}
When $\Phi_{\mathrm{biased}}$ is minimised to a value below a threshold
$\varepsilon_{\mathrm{tol}}$, both conditions are jointly satisfied:
\begin{enumerate}
    \item The stealthy constraint $S(\mathbf{k})\approx 0$ holds for all
          $\mathbf{k}\in\mathcal{K}$ (from $\Phi_{\mathrm{CC}}\approx 0$).
    \item All particles lie outside the sphere of radius $R_f$ centred at
          $\mathbf{r}_c$ (from $\phi_{\mathrm{ex}}\approx 0$, which
          requires $r_{ic}\ge R_f$ for every particle).
\end{enumerate}

% ------------------------------------------------------------------
\section{Force from the exclusion field}
% ------------------------------------------------------------------

The force on particle $i$ from the exclusion field is
\begin{equation}
    \mathbf{F}_i^{\mathrm{ex}}
    = -\nabla_{\mathbf{r}_i}\,\phi_{\mathrm{ex}}
    = \begin{cases}
        \displaystyle
        \frac{R_f}{r_{ic}^{3}}\,(\mathbf{r}_i - \mathbf{r}_c)
        & \text{if } r_{ic} < R_f,\\[8pt]
        \mathbf{0} & \text{otherwise.}
      \end{cases}
    \label{eq:force}
\end{equation}
This force points \emph{radially outward} from $\mathbf{r}_c$ and grows
without bound as $r_{ic}\to 0$, acting as a hard barrier that prevents
any particle from entering the probe sphere.

The total force on particle $i$ during minimisation is
\begin{equation}
    \mathbf{F}_i
    = \mathbf{F}_i^{\mathrm{CC}} + \mathbf{F}_i^{\mathrm{ex}},
\end{equation}
where $\mathbf{F}_i^{\mathrm{CC}} = -\nabla_{\mathbf{r}_i}\Phi_{\mathrm{CC}}$
is the standard stealthy force already present in the codebase.

% ------------------------------------------------------------------
\section{Feasibility scan and maximum hole radius}
% ------------------------------------------------------------------

The probe radius $R_f$ is treated as a fixed external parameter that is
stepped upward from $R_{\min}$ to $R_{\max}$ in increments $\Delta R$:
\begin{equation}
    R_f \;\in\; \{R_{\min},\; R_{\min}+\Delta R,\; R_{\min}+2\Delta R,\;\ldots\}.
\end{equation}
At each value of $R_f$:
\begin{enumerate}
    \item $N_{\mathrm{trial}}$ independent random initial configurations
          are generated and minimised under $\Phi_{\mathrm{biased}}$.
    \item If \emph{any} minimisation reaches
          $\Phi_{\mathrm{biased}} < \varepsilon_{\mathrm{tol}}$,
          the void of radius $R_f$ is \textbf{feasible}: record
          $R_c = R_f$ and advance to the next step.
    \item If \emph{all} $N_{\mathrm{trial}}$ attempts fail, the void is
          \textbf{infeasible}: stop the scan.
\end{enumerate}
The output is the maximum feasible hole radius $R_c$ and the
corresponding stealthy configuration.

% ------------------------------------------------------------------
\section{Dimensionless output}
% ------------------------------------------------------------------

Results are reported in two units:
\begin{itemize}
    \item $R_c$ in absolute (simulation) length units.
    \item $R_c \cdot k_{\max}$ (dimensionless), where
          $k_{\max}$ is the outer edge of the outermost excluded shell:
          \begin{equation}
              k_{\max} = \max_{n}\bigl(k_1^{(n)} + \delta^{(n)}\bigr).
          \end{equation}
\end{itemize}
Running the scan across many independent seeds (SLURM array jobs, each
with a different random initial condition) samples the distribution of
$R_c$ over the stealthy ground-state ensemble, giving the hole-size
distribution $p(R_c)$.

\end{document}
